\section{Experiment}
\label{subsec:experiment}

\paragraph*{Dataset}
We evaluate our method on the RealX3D low-light benchmark dataset~\cite{liu2025realx3d}.
RealX3D contains nine real-world scenes captured under low-light conditions,
which features complex scene geometry, denser multi-view capture, wider viewpoint
diversity, and spatially non-uniform light setup. Since the BlueHawaii scene was
used as the validation scene in the 3DRR challenge, we follow this protocol and
conduct all comparisons on the remaining eight scenes. For each scene, the
training set provides 27--32 low-light sRGB images together with the corresponding
calibrated camera poses, and the task is to synthesize 5--6 novel-view
images under normal illumination.

\paragraph*{Comparison Methods}
We compare our method with representative low-light 3D reconstruction approaches
that operate directly on sRGB inputs, including vanilla 3DGS~\cite{kerbl20233d},
AlethNeRF~\cite{cui2024aleth}, I2NeRF~\cite{liu2025i2nerf}, LITA-GS~\cite{zhou2025lita},
and Luminance-GS~\cite{cui2025luminance}, as summarized in Table~\ref{tab:comparison_main}. These
baselines cover both NeRF-based methods (AlethNeRF and I2NeRF) and 3DGS-based
methods (vanilla 3DGS, LITA-GS, and Luminance-GS). For fair comparison, neither
our method nor the compared methods relies on separate pre-processing exposure
enhancement or post-processing image restoration modules.

\paragraph*{Implementation Details}
We extract monocular depth, offline structure priors, and fixed-pose COLMAP sparse
points for each scene, and use the same Simple3DGS backbone with spherical harmonics
of degree 3 in all stages. Low-light training images are converted into enhanced
supervision on the fly, while Stage 3 additionally uses an adaptive proxy target.
Experiments were conducted on a cloud-hosted NVIDIA GPU with 32 GB VRAM.

\begin{table*}[!t]
\centering
\scriptsize
\setlength{\tabcolsep}{4.0pt}
\renewcommand{\arraystretch}{1.10}
\resizebox{\textwidth}{!}{%
\begin{tabular}{@{}lccccccccc@{}}
\toprule
Method & Chocolate & Cupcake & GearWorks & Laboratory
& MilkCookie & Popcorn & Sculpture & Ujikintoki & Average \\
\midrule
& \multicolumn{9}{c}{PSNR$\uparrow$} \\
\cmidrule(lr){2-10}
3DGS & 7.89 & 4.42 & 6.90 & 7.59 & 6.15 & 7.46 & 5.85 & 6.23 & 6.56 \\
AlethNeRF & 11.53 & \thirdcell{13.49} & 8.82 & 14.14 & \thirdcell{13.49}
& 13.25 & \thirdcell{12.22} & 14.12 & 12.63 \\
I2NeRF & \bestcell{19.77} & 12.68 & \secondcell{12.22}
& \thirdcell{16.34} & \secondcell{14.78} & \thirdcell{17.08}
& 9.64 & \secondcell{19.02} & \thirdcell{15.19} \\
LITA-GS & \thirdcell{17.94} & 13.07 & \thirdcell{10.90}
& \secondcell{17.56} & 12.74 & \bestcell{18.97}
& \secondcell{12.84} & \thirdcell{18.82} & \secondcell{15.36} \\
Luminance-GS & 7.33 & \secondcell{14.82} & 8.09 & 8.88 & 9.67
& 11.33 & 9.79 & 9.36 & 9.91 \\
\textbf{Ours} & \secondcell{19.09} & \bestcell{15.57} & \bestcell{15.61}
& \bestcell{22.63} & \bestcell{16.01} & \secondcell{18.82}
& \bestcell{15.52} & \bestcell{19.81} & \bestcell{17.88} \\
\midrule
& \multicolumn{9}{c}{SSIM$\uparrow$} \\
\cmidrule(lr){2-10}
3DGS & 0.084 & 0.064 & 0.081 & 0.041 & 0.052 & 0.062 & 0.018
& 0.062 & 0.058 \\
AlethNeRF & 0.406 & 0.552 & 0.217 & 0.484 & 0.567 & 0.452 & 0.208
& 0.544 & 0.429 \\
I2NeRF & \bestcell{0.571} & \thirdcell{0.608} & \secondcell{0.511}
& \secondcell{0.630} & \secondcell{0.649} & \thirdcell{0.543}
& 0.277 & \secondcell{0.666} & \secondcell{0.557} \\
LITA-GS & \secondcell{0.541} & \secondcell{0.643} & 0.344
& \thirdcell{0.597} & \thirdcell{0.573} & \secondcell{0.557}
& \secondcell{0.343} & \thirdcell{0.656} & \thirdcell{0.532} \\
Luminance-GS & 0.362 & 0.536 & \thirdcell{0.411} & 0.372 & 0.532
& 0.462 & \thirdcell{0.293} & 0.463 & 0.429 \\
\textbf{Ours} & \thirdcell{0.534} & \bestcell{0.806} & \bestcell{0.647}
& \bestcell{0.711} & \bestcell{0.740} & \bestcell{0.665}
& \bestcell{0.482} & \bestcell{0.766} & \bestcell{0.669} \\
\midrule
& \multicolumn{9}{c}{LPIPS$\downarrow$} \\
\cmidrule(lr){2-10}
3DGS & 0.645 & 0.625 & 0.658 & 0.610 & 0.653 & 0.650 & 0.733
& 0.683 & 0.657 \\
AlethNeRF & 0.692 & 0.612 & 0.671 & 0.696 & 0.725 & 0.688 & 0.891
& 0.646 & 0.703 \\
I2NeRF & \thirdcell{0.561} & \thirdcell{0.493} & \thirdcell{0.543}
& \thirdcell{0.486} & \thirdcell{0.544} & \thirdcell{0.488}
& \thirdcell{0.648} & \bestcell{0.474} & \thirdcell{0.530} \\
LITA-GS & \secondcell{0.557} & \secondcell{0.326} & \secondcell{0.521}
& \secondcell{0.435} & \secondcell{0.478} & \secondcell{0.448}
& \secondcell{0.588} & \thirdcell{0.497} & \secondcell{0.481} \\
Luminance-GS & 0.779 & 0.534 & 0.737 & 0.614 & 0.766 & 0.625 & 0.655
& 0.882 & 0.699 \\
\textbf{Ours} & \bestcell{0.517} & \bestcell{0.303} & \bestcell{0.460}
& \bestcell{0.392} & \bestcell{0.424} & \bestcell{0.444}
& \bestcell{0.492} & \secondcell{0.478} & \bestcell{0.439} \\
\bottomrule
\end{tabular}%
}
\caption{Quantitative comparison with other baseline methods on 8 scenes
from RealX3D lowlight dataset, including \textit{Chocolate}, \textit{Cupcake}, \textit{GearWorks},
\textit{Laboratory}, \textit{MilkCookie}, \textit{Popcorn}, \textit{Sculpture}, \textit{Ujikintoki}. We reported three
rendering metrics (PSNR, SSIM, LPIPS). The first, second and third results
are highlighted in dark, medium, and light orange, respectively.}
\label{tab:comparison_main}
\end{table*}


Our final pipeline contains three stages. The first stage 
runs for 15K iterations, initializes Gaussians from the filtered COLMAP sparse cloud
plus 5000 random points, and is the only stage with densification. The second stage
runs for 5K iterations without densification and enables
sparse-guided regularization with 4096 sampled active Gaussians, $K{=}3$, and weight
$0.1$. The final stage also
runs for 5K iterations, warm-starts from the geometry base, freezes the main geometry, and
performs YCbCr-space appearance refinement with $\lambda_Y{=}0.5$,
$\lambda_{CbCr}{=}0.16$, $\lambda_{chr}{=}5{\times}10^{-4}$, and both luminance and
chroma weight maps enabled.

\paragraph*{Experimental Results}
We compare our method with 3DGS, AlethNeRF, I2NeRF, LITA-GS, and
Luminance-GS on the eight evaluation scenes excluding BlueHawaii.
Table~\ref{tab:comparison_main} reports PSNR, SSIM, and LPIPS. Our method achieves the
best average performance on all three metrics, reaching 17.88 PSNR, 0.669
SSIM, and 0.439 LPIPS. It ranks first on 6 out of 8 scenes in PSNR, 7 out
of 8 scenes in SSIM, and 7 out of 8 scenes in LPIPS. Compared with the
strongest 3DGS-based baseline, LITA-GS, our final model improves the scene
average by 2.52 dB in PSNR and reduces LPIPS from 0.481 to 0.439. The gains
are particularly clear on challenging scenes such as \textit{GearWorks}, \textit{Laboratory},
and \textit{Sculpture}, where reliable geometry and decoupled appearance refinement
are both important. 

Figure~\ref{fig:qualitative_results} shows representative visual comparisons
on five scenes. Existing methods often suffer from unstable enhancement,
severe texture loss, or distorted color restoration under severe low-light
conditions. NeRF-based methods, like AlethNeRF and I2NeRF, produce messy
results with excessive scene noise and chaotic colors when given extremely
dark input images. Luminance-GS tends to over-enhance low-light scenes,
leading to haze-like appearance renderings and noticeable loss of contrast
rather than natural brightness recovery. LITA-GS exhibits strong sensitivity
to varying scenes. Specifically, it leads to excessive enhancement of brightness
and contrast in certain scenarios, such as \textit{Popcorn} and
\textit{Chocolate}, while in other scenarios like \textit{GearWorks}, it
suffers from insufficient enhancement and texture loss. In contrast, our
method consistently produces clearer scene boundaries, more plausible
brightness, and more natural color appearance, avoiding the strong color cast
or washed-out appearance observed in several baselines. The improvements are
especially visible in the cluttered \textit{Chocolate}, \textit{Laboratory}
and \textit{Popcorn} scenes and in the geometry-sensitive \textit{Sculpture}
scene, where our results preserve fine structures.

\paragraph*{Ablation Study}
Table~\ref{tab:stage_ablation} shows a clear separation between the roles of
the three stages. Stage~1 already gives a reasonable starting point, while
adding Stage~2 changes the image metrics only slightly on average
(14.74$\rightarrow$14.84 PSNR, 0.598$\rightarrow$0.599 SSIM, and
0.465$\rightarrow$0.463 LPIPS). This behavior is consistent with the goal of
Stage~2: it mainly stabilizes geometry under no-densify refinement instead of
directly optimizing perceptual appearance. The major improvement appears after
Stage~3, where the average performance rises to 17.88 PSNR, 0.669 SSIM, and
0.439 LPIPS, with gains on all scenes. This indicates that once geometry is
sufficiently stabilized, the remaining bottleneck lies in appearance refinement
under low-light conditions, which is effectively handled by the geometry-frozen
decoupled refinement stage.

\begin{table*}[!t]
\centering
\scriptsize
\setlength{\tabcolsep}{3.8pt}
\renewcommand{\arraystretch}{1.10}
\resizebox{\textwidth}{!}{%
\begin{tabular}{@{}lccccccccc@{}}
\toprule
Stage & Chocolate & Cupcake & GearWorks & Laboratory
& MilkCookie & Popcorn & Sculpture & Ujikintoki & Average \\
\midrule
& \multicolumn{9}{c}{PSNR$\uparrow$} \\
\cmidrule(lr){2-10}
Stage 1 only & \secondcell{16.84} & \secondcell{12.49} & \thirdcell{14.25}
& \thirdcell{15.64} & \thirdcell{13.91} & \thirdcell{15.68}
& \secondcell{13.47} & \thirdcell{15.63} & \thirdcell{14.74} \\
Stages 1+2 & \thirdcell{16.72} & \thirdcell{12.46} & \secondcell{14.51}
& \secondcell{15.81} & \secondcell{14.10} & \secondcell{15.94}
& \thirdcell{13.43} & \secondcell{15.76} & \secondcell{14.84} \\
Stages 1+2+3 & \bestcell{19.09} & \bestcell{15.57} & \bestcell{15.61}
& \bestcell{22.63} & \bestcell{16.01} & \bestcell{18.82}
& \bestcell{15.52} & \bestcell{19.81} & \bestcell{17.88} \\
\midrule
& \multicolumn{9}{c}{SSIM$\uparrow$} \\
\cmidrule(lr){2-10}
Stage 1 only & \thirdcell{0.4922} & \secondcell{0.7430}
& \thirdcell{0.5821} & \thirdcell{0.6024} & \secondcell{0.6754}
& \thirdcell{0.5781} & \secondcell{0.4471} & \thirdcell{0.6635}
& \thirdcell{0.5980} \\
Stages 1+2 & \secondcell{0.4940} & \thirdcell{0.7388}
& \secondcell{0.5853} & \secondcell{0.6077} & \thirdcell{0.6742}
& \secondcell{0.5816} & \thirdcell{0.4448} & \secondcell{0.6666}
& \secondcell{0.5991} \\
Stages 1+2+3 & \bestcell{0.5341} & \bestcell{0.8060} & \bestcell{0.6468}
& \bestcell{0.7112} & \bestcell{0.7404} & \bestcell{0.6649}
& \bestcell{0.4820} & \bestcell{0.7662} & \bestcell{0.6690} \\
\midrule
& \multicolumn{9}{c}{LPIPS$\downarrow$} \\
\cmidrule(lr){2-10}
Stage 1 only & \thirdcell{0.5436} & \secondcell{0.3143}
& \thirdcell{0.5063} & \thirdcell{0.4067} & \thirdcell{0.4788}
& \secondcell{0.4795} & \thirdcell{0.5007} & \secondcell{0.4917}
& \thirdcell{0.4652} \\
Stages 1+2 & \secondcell{0.5415} & \thirdcell{0.3155}
& \secondcell{0.5031} & \secondcell{0.4006} & \secondcell{0.4749}
& \thirdcell{0.4798} & \secondcell{0.4966} & \thirdcell{0.4938}
& \secondcell{0.4632} \\
Stages 1+2+3 & \bestcell{0.5168} & \bestcell{0.3032} & \bestcell{0.4597}
& \bestcell{0.3924} & \bestcell{0.4243} & \bestcell{0.4442}
& \bestcell{0.4921} & \bestcell{0.4777} & \bestcell{0.4388} \\
\bottomrule
\end{tabular}%
}
\caption{Three-stage ablation of the proposed pipeline. Stage 1
performs geometry bootstrap, Stage 2 refines geometry without
densification, and Stage 3 conducts geometry-frozen appearance
refinement. Better results are highlighted in darker orange.}
\label{tab:stage_ablation}
\end{table*}


The effect of sparse-guided regularization is better reflected by
geometry-oriented metrics than by image metrics. As shown in
Table~\ref{tab:sparse_geometry_ablation}, enabling this term consistently
reduces both the average support distance and the Chamfer-like mean, while
improving G$\rightarrow$S Coverage@0.10. Although the absolute gains are
modest, they are consistent across all three measures, suggesting that the
module works as intended: it gently pulls active Gaussians toward more
reliable sparse support without destabilizing the rendering optimization.

\par\medskip
\noindent\begin{minipage}{\columnwidth}
\centering
\scriptsize
\setlength{\tabcolsep}{3.5pt}
\renewcommand{\arraystretch}{1.10}
\resizebox{\columnwidth}{!}{%
\begin{tabular}{@{}lccc@{}}
\toprule
Setting & Support Dist.$\downarrow$ & Chamfer Mean$\downarrow$
& G$\rightarrow$S Cov.@0.10$\uparrow$ \\
\midrule
w/o sparse-guided & \secondcell{0.3736} & \secondcell{0.3722}
& \secondcell{0.6121} \\
w/ sparse-guided & \bestcell{0.3712} & \bestcell{0.3696}
& \bestcell{0.6164} \\
\bottomrule
\end{tabular}
}
\captionof{table}{Scene-averaged geometry ablation of sparse-guided regularization
in Stage 2. Metrics are computed on the final Stage-2 checkpoints using
active Gaussians against COLMAP sparse points.}
\label{tab:sparse_geometry_ablation}
\end{minipage}\par\medskip


Table~\ref{tab:stage3_weightmap_ablation} evaluates the key components of the
final appearance refinement stage. The full model performs best overall.
Removing the CbCr weight map leads to a small but consistent degradation,
showing that branch-specific chroma weighting is helpful but not the dominant
factor. In contrast, removing the Y weight map causes a clearly larger drop,
especially in PSNR (17.88$\rightarrow$17.64), indicating that luminance-aware
reweighting contributes more directly to restoring low-light visibility. The
largest degradation comes from removing the chroma residual branch, which
reduces performance to 17.27 PSNR, 0.6486 SSIM, and 0.4505 LPIPS. This result
shows that brightness correction alone is insufficient: explicit chroma
adjustment is still necessary to suppress color cast and recover natural
appearance.
\noindent\begin{minipage}{\columnwidth}
\centering
\small
\setlength{\tabcolsep}{6pt}
\renewcommand{\arraystretch}{1.10}
\begin{tabular}{@{}lccc@{}}
\toprule
Setting & PSNR$\uparrow$ & SSIM$\uparrow$ & LPIPS$\downarrow$ \\
\midrule
Full model & \bestcell{17.88} & \bestcell{0.6690}
& \bestcell{0.4388} \\
w/o CbCr weight map & \secondcell{17.87} & \secondcell{0.6680}
& \secondcell{0.4397} \\
w/o Y weight map & \thirdcell{17.64} & \thirdcell{0.6679}
& \thirdcell{0.4400} \\
w/o chroma residual branch & 17.27 & 0.6486 & 0.4505 \\
\bottomrule
\end{tabular}
\captionof{table}{Scene-averaged ablation of the Stage 3 design. Metrics are
averaged over the eight evaluated scenes.}
\label{tab:stage3_weightmap_ablation}
\end{minipage}

